\documentclass{article}
\usepackage[left=2cm, right=2cm, top=2cm, bottom=2cm]{geometry}
\usepackage{graphicx}
\usepackage{amsmath}
\usepackage{amssymb}
\usepackage{amsthm}
\usepackage{fancyhdr}
\usepackage{verbatim}
\usepackage{listings}
\usepackage{xcolor}
\usepackage{pdfpages}
\usepackage{cancel}
\usepackage{physics}
\usepackage{esint}

\lstset{
    language=C++,
    basicstyle=\ttfamily\footnotesize,
    keywordstyle=\color{blue},
    commentstyle=\color{green},
    stringstyle=\color{red},
    numbers=left,
    numberstyle=\tiny\color{gray},
    frame=single,
    breaklines=true,
    captionpos=b,
}


\title{MTH 553: Lecture \# 28}
\author{Cliff Sun}

\newtheorem{theorem}{Theorem}[section]
\newtheorem{lemma}[theorem]{Lemma}
\newtheorem{definition}[theorem]{Definition}
\newtheorem{conjecture}[theorem]{Conjecture}
\newtheorem{proposition}[theorem]{Proposition}
\newtheorem{corollary}[theorem]{Corollary}
\newtheorem{one minute paper}[theorem]{One Minute Paper}

\pagestyle{fancy}
\lhead{\textbf{\thepage}\ \ \nouppercase{\rightmark}}
\chead{MTH 553: Lecture \# 28}
\rhead{Cliff Sun}

\begin{document}

\maketitle

\section*{Lecture Span}
\begin{enumerate}
    \item Poissons equation
\end{enumerate}

\section*{Continuous $L^{\infty}$ dependence}

\begin{equation}
    \begin{cases}
        \triangle u_1 = f \\
        u_1 = g_1
    \end{cases}
\end{equation}

\begin{equation}
    \begin{cases}
        \triangle u_2 = f \\
        u_2 = g_2
    \end{cases}
\end{equation}

With $|g_1 - g_2| < \epsilon$ on $\partial \Omega$ implies that $|u_1 - u_2| \leq \epsilon$ on $\Omega$. 

Proof, use the maximum principle of $u$ on $\partial \Omega$ (which is where $g$ is defined anyways). 

\section*{Distributions}

Let $\Omega$ be an open set on $\mathbb{R}^n$. 

\begin{definition}
   (Convergence in $C_0^{\infty}$) Let $v_{j}, v \in C_{0}^{\infty}(\Omega)$. Say $v_j \rightarrow v$ in $C_0^{\infty}$ if 
   as $j \rightarrow \infty$ if 
   \begin{enumerate}
    \item Exists compact set $k \subset \Omega$ such that $v_j$ are supported in $K$. 
    \item For all $\alpha$ (multi-indices), we have that $D^{\alpha}v_j \rightarrow D^{\alpha}v$ uniformly in $\Omega$ as $j \rightarrow \infty$.
   \end{enumerate}
\end{definition}

For example, $$v + \frac{1}{j}v \rightarrow v$$

\begin{definition}
    \begin{enumerate}
        \item (Distributions) Call a linear mapping 
        \begin{equation}
            F: C_{0}^{\infty}(\Omega) \rightarrow \mathbb{R}
        \end{equation}
        A \underline{distribution} if it is continuous, meaning $F(v_j) \rightarrow F(V)$ as $v_j \rightarrow v$. 
        \item (Convergence of distributions) Let $F_k, F$ be distributions. Say $F_k \rightarrow F$ if 
        \begin{equation}
            F_k(v) \rightarrow F(v)    
        \end{equation}
        As $k \rightarrow \infty$ whenever $v$ is a test function and is in $C_{0}^{\infty}(\Omega)$. 
    \end{enumerate}
\end{definition}

Notations, $\langle F, V \rangle$ or $\int_{\Omega}Fvdx$ for $F(v)$. Write $D = C_{0}^{\infty}(\Omega)$ and $D'$ is the space of distributions or the dual space of $D$. 

\subsection*{Examples}

All integrable functions are contained in all signed measures which is also contained within distributions. Given $f \in L^1(\Omega)$, define a distribution 

\begin{equation}
    F_f(v) = \int_{\Omega}fv dx
\end{equation}

Check that $F_f(v_j) \rightarrow F_f(v)$ as $v_j \rightarrow v$. This example also motivates the notation 

\begin{equation}
    \int_{\Omega}Fv dx
\end{equation}

for a distribution. Given a signed Borel measure $\mu$, we may define a distribution as 

\begin{equation}
    F_{\mu}(v) = \int_{\Omega}vd\mu
\end{equation}

Whatever $\mu$ means. Fix $y \in \mathbb{R}^n$. Define a distribution $\delta_y$ by: 

\begin{equation}
    \delta_y(v) = v(y)
\end{equation}

This is also called the delta distribution concentrated at $y$. Notationally, 

\begin{equation}
    \int_\Omega\delta_y v dx = v(y)
\end{equation}

We can also approximate the delta distribution. Let $f \in  L^1(\mathbb{R}^n)$ with 

\begin{equation}
    \int f dx = 1
\end{equation}

For $\epsilon > 0$, rescale $f$ as follows: Letting $$f_{\epsilon}(x) = \frac{1}{\epsilon^{n}}f\left( \frac{x}{\epsilon} \right)$$

Note that 

\begin{equation}
    \int f_{\epsilon}dx = 1
\end{equation}

And 

\begin{equation}
    F_{f_{\epsilon}} \rightarrow \delta_0, \quad \epsilon \rightarrow 0
\end{equation}

We prove both statements:

\begin{proof}
    For all $v \in C_0^{\infty}$: 
\end{proof}

Then, we check 

\begin{align*}
    \langle F_{f_\epsilon}, v\rangle &= \int\frac{1}{\epsilon^n}f\left(\frac{x}{\epsilon}\right)v(x)dx \\
    &= \int_{\mathbb{R}^n}f(y)v(\epsilon y)dy\\
    &= \int_{\mathbb{R}^n}f(y)v(0) dy, \quad \text{by dominated convergence in measure theory} \\
    &= v(0)\int_{\mathbb{R}^n}f(y)dy \\
    &= v(0) = \langle \delta_0, v \rangle
\end{align*}

\begin{definition}
    Derivatives of distributions. Given a distribution $u$ and a linear differential operator $L$, then define a new distribution
    $$Lu(v) \equiv u(L' v)$$ where $L'$ is the adjoint of $L$.
\end{definition}

Notationally, 

\begin{equation}
    \langle Lu, v \rangle \equiv \langle u, L' v \rangle
\end{equation}

For instance, $L = \partial_{x_1}$, then 

\begin{equation}
    \langle \partial_{x_1} u , v \rangle = -\langle u, \partial_{x_1}v \rangle
\end{equation}

Therefore, every distribution is infinitely differentiable. Note, to say that $Lu = f$ in terms of distributions that 

\begin{equation}
    \langle Lu, v \rangle = \langle f, v \rangle, \quad \forall v \in C_0^{\infty}\left( \Omega \right)
\end{equation}
\begin{equation}
    \langle u, L' v \rangle = \langle f, v \rangle
\end{equation}

If $u,f \in L^1(\Omega)$ (are integrable), then 

\begin{equation}
    \int u L' v dx = \int fv dx
\end{equation} 

If $Lu = f$ weakly where $u, f \in L^1(\Omega)$, then 

\begin{equation}
    Lu = f
\end{equation}

as distributions. 

\end{document}